\section{Introduction}
\label{sec:intro}
Shapley-value attributions \citep{lundberg2017unified} are used to justify individual predictions of tabular models in credit, health and similar settings. Outside tree models and linear models, where exact algorithms exist \citep{lundberg2020local}, the attribution is \emph{estimated} by sampling coalitions or permutations, with a budget of model evaluations that the practitioner fixes. A common practice is to run the explainer once for an instance and report the resulting vector as the reason for that prediction. If the estimator is noisy, a second run could report a different top feature, a different sign or a different order, and the reported ``reason'' would then be partly an artefact of the random seed.

Prior work quantifies how estimator variance falls with the number of samples \citep{covert2021improving,mitchell2022sampling,musco2025leverage}, how explanations differ between methods \citep{krishna2024disagreement}, and how the background set matters \citep{yuan2023background,chen2020true}. What a practitioner needs is different: for a trained model, a given instance and a realistic budget, \emph{how far can one run be treated as the attribution of that prediction, and where does that fail?} Answering this requires deciding what ``reliable'' means, against which reference, and under which conditions; a single agreement number between two runs is not enough, because it cannot tell noise from bias and cannot say whether the noise matters for what a reader would conclude.

We study this question empirically with four separable notions of reliability (Section~\ref{sec:method}): repeatability between independent runs, accuracy against the exact Shapley vector for the same model, background and value function, resolvability of adjacent features against the run-to-run noise, and the size of explainer noise relative to other variation that a deployed explanation already carries. We evaluate them in the regime where sampling noise exists at all: Kernel SHAP enumerates every coalition once the number of samples reaches $2^d-2$, so for $d\le10$ the library default is deterministic; our core datasets therefore have $d=20$ and $d=16$, and two lower-dimensional datasets serve as controls.

\paragraph{Findings (scope: $d\le20$, background of 50 rows, one fit per model).} (i) With the library's default sample count and $\ell_1$ selection switched off, a single Kernel SHAP run recovers the top-3 features with mean hit rate between $\nGerOneTopHitEstMin$ and $\nGerOneTopHitEstMax$ on German Credit and $\nAcsOneTopHitEstMin$ and $\nAcsOneTopHitEstMax$ on ACSEmployment, and the decision rule classes all eight cells as reliable. Two of the four models are not informative about sampling noise (LR is additive, so the estimate is exact; XGBoost and RF have an exact algorithm that one would use instead), so we also state the MLP alone: hit rate $\nGerMlpOneTopHitEst$ (20 instances) and $\nAcsMlpOneTopHitEst$. Permutation SHAP behaves similarly. (ii) The library's default $\ell_1$ setting is the one configuration in which the rule is not met (``conditional'') on German Credit; the error is a deterministic bias that repeating runs does not reduce. This is a configuration pitfall of one implementation. (iii) Reliability is not uniform inside an attribution vector: a single run reproduces the order of an adjacent pair with probability $\nGerPOrderMin$--$\nGerPOrderMax$ on German Credit, and $\nGerPairsResMin$--$\nGerPairsResMax$ of adjacent pairs are separated by twice the run-to-run noise, with the unresolved pairs at the lower ranks. (iv) At a background of 50 rows, redrawing the background changes the attribution vector with variance $\nGerBgRatioMin$--$\nGerBgRatioMax$ (German Credit) and $\nAcsBgRatioMin$--$\nAcsBgRatioMax$ (ACSEmployment) times that of the explainer seed; since a redraw changes the target rather than the estimate, this describes sizes of variation, not accuracy. Section~\ref{sec:discussion} states what these results do not cover.

\paragraph{Contributions.} (1) An evaluation design that combines standard stability and error measures (and a decision rule fixed in the study plan before the runs) with exact targets, for the single-run practice at the library's default budget; the metrics themselves are not new. (2) A controlled evaluation on two stochastic and two deterministic-regime datasets, four model families and two sampling explainers, with nested-bootstrap uncertainty and exact targets for every explained model. (3) A measured configuration pitfall: the default $\ell_1$ setting of \texttt{shap} 0.51.0's \texttt{KernelExplainer} biases single-run attributions at the default budget on a $d=20$ dataset (the arm was run at $1\times$ only), and averaging runs at equal total budget does not improve on one larger run in our data (a small effect). (4) All code and tidy results.
